\section{Interpretation of the generated nonlinear terms}
\label{sec:interpretation}

The worked example makes the origin of each monomial family transparent.

\subsection{The SII family}

A term $S_iI_\ell I_j$ appears when the susceptible factor in the infection
observable $S_iI_j$ is differentiated. One infectious factor, $I_j$, was
already present in the observable; the second, $I_\ell$, comes from an incoming
source of city $i$. These terms therefore encode pairs of infectious factors
sharing the same receiving susceptible population.

\subsection{The SIE and SEE families}

A term $S_iI_\ell E_j$ appears when the susceptible factor in $S_iE_j$ is
differentiated. Linear closure then replaces infectious factors by their
exposed counterparts and produces $SEE$ coordinates. These coordinates are
required because exposed and infectious variables are coupled by the linear
compartmental progression.

\subsection{The SSI and SSE families}

A term $S_iS_jI_\ell$ appears when the exposed factor in $S_iE_j$ is
differentiated through the positive transmission contribution to $\dot E_j$.
It links the susceptible factor already present at receiving city $i$ to the
susceptible population and incoming infectious source structure of city $j$.
Linear closure converts the infectious factor into an exposed factor and adds
the corresponding $SSE$ coordinate.

\subsection{Effect of diagonal channels}

Diagonal channels create terms such as
\[
S_iI_i^2,
\qquad
S_iE_iI_i,
\qquad
S_i^2I_i,
\]
and their linearly closed variants. These are not algebraic anomalies. They
represent aggregate within-city transmission and naturally contain repeated
city indices and repeated scalar factors.

\subsection{What the example does and does not establish}

The example establishes the exact operation of Algorithm~1 on a concrete
network and verifies the reported dictionary dimensions. It illustrates how
network support controls which polynomial interactions are generated. It does
not claim that the resulting dictionary is the smallest possible invariant or
approximately invariant dictionary, nor does it provide a trajectory-error
bound for a finite truncation.
